\subsection{训练证据如何转化为下一代合同}
当前 \code{engine/evolution.py} 将方向诊断、条件覆盖和信号组合编码为结构化证据。\code{training\_stop} 使用 $\max(0,\min(800,\lfloor0.55T\rfloor)-\max(15,h+1))$ 作为训练前缀终点；其后保留 purge 间隔再计算后续描述性诊断。选择符号来自训练均值，后续统计不倒灌为该轮方向选择。该拆分仍处于 A 段，不能等同于 B/H 独立确认。

子代动作包括 condition（完整条件路径）、reverse（冻结反向假设）、interaction（跨构件有界交互），以及 horizontal、cross\_family、forest、crossover 与记忆 probe。失败或未知父结构仍可能包含有用的训练异质性线索；这只允许产生新的探索假设，不自动授予有效性。子代记录 parent、donor、operator 与 depth，生成表达式后仍须通过 DSL、机制审查和新冻结。最大深度为 2，队列与模型调用均受预算约束。

\subsection{与无限制自修改的区别}
此处 RSI 的递归对象是\important{研究问题、候选结构、训练条件和搜索分配}。模型不会重写 Python 评估器，也不会自动训练或更新自身权重。执行阶段的数字、断言结果与未完成状态由确定性代码产生；模型负责提出与组织研究解释。这使“下一代研究更具体”成为可审计行为，但是否带来更高质量因子仍需要持续实验验证。

\subsection{组合 Gate 的执行含义}
多结构组合并非将表达式简单相加。系统保留各结构的原覆盖条件，在持仓路径层建立联合规则：等资本并联、冲突区域空仓、至少两个独立构造同向。重复构造需要合并，未知收益不能填零。研究回报代理与 RQAlpha 执行核验分开记录；缺少执行价格或必要确认时，控制台显示失败或等待，不给出虚构净值。

\subsection{性能优化的层次与可测量边界}
\begin{enumerate}[leftmargin=2em]
\item \textbf{共享缓存：} 冻结配置经过规范化后计算缓存身份，父进程和 spawn 子进程使用一致表示，减少重复构建面板。
\item \textbf{切分并行：} 独立进程按日期读取临时 Parquet，内存预算限制 worker 数；原生线程限制为 1，避免 BLAS 与进程池相乘。
\item \textbf{初始网格并行：} Linux 专用 \code{engine.initial\_search} 对不自适应的初始坐标分片，CPU affinity 与共享项目锁约束资源和审计；它不同于递归队列，不能直接等同于完整 RSI 并行运行。
\item \textbf{GPU nuisance 回归：} 可选 CuPy 路径流式累计充分统计量并限制稀疏工作区。它加速回归子步骤，不代表整个森林或全流程均在 GPU 上运行；GPU 依赖与可用显存需要单独配置。
\item \textbf{产品读取：} 前端请求结束后再等待 5 秒轮询，切换批次取消旧请求；默认二维地图，三维按需加载，减少展示对研究资源的占用。
\end{enumerate}
正式、工程和 fast 模式的计算预算不同，不能将 fast 耗时与 formal 精度作无条件等价比较。历史阶段耗时需在同机、同数据、同冻结配置下复核；本次报告更新不额外声称全市场最新加速比。
\subsection{控制指令与研究状态分离}
控制台的 pause/stop 回执表示请求已原子落盘，研究进程在检查点边界生效后才改变实际状态。每批次控制锁串行化操作；重复 resume/restart 不再对存活的已登记 PID 启动第二个恢复进程。已完成批次拒绝继续修改运行状态。

创建恢复进程前，服务先核对代码哈希和 Python/依赖环境，完整数据、模型和产物身份由引擎继续核验。缺少 PID 的外部 CLI 批次不被自动判定为退出。该机制是本机可靠性控制，不是多用户认证或远程运维权限系统；商业部署仍需身份认证、租户隔离、访问策略及受管理的作业服务。